\begin{ExerciseBlock}
\ExerciseHeading{}{Verständnisfragen}{angular_momentum_checks}
Entscheiden Sie, welche Aussagen wahr oder falsch sind, und begründen Sie Ihre Entscheidung.
\begin{enumerate}
\item Die Einheit des Drehimpulses ist Ns/m.\TrueFalseAnswer{falsch}
\item Bezogen auf einen festen Punkt im Inertialsystem entspricht das resultierende Drehmoment $\vec M$ der zeitlichen Änderung des Drehimpulses $\vec L$.\TrueFalseAnswer{wahr}
\item Der Wert des Drehimpulses $\vec{L}$ hängt nicht von der Wahl des Referenzpunktes ab.\TrueFalseAnswer{falsch}
\item Drehimpulserhaltung gilt unabhängig von der Wahl des Bezugspunkts, auch wenn das Drehmoment bezüglich eines der Punkte von null verschieden ist.\TrueFalseAnswer{falsch}
\item Zum Drehimpuls tragen nur Geschwindigkeiten senkrecht zum Verbindungsvektor $\vec{r}$ bei.\TrueFalseAnswer{wahr}
\item Der Drehimpulsvektor $\vec{L}$ liegt in der Drehebene einer Rotationsbewegung.\TrueFalseAnswer{falsch}
\item Bei einer Drehung um eine feste Achse beschreibt das Massenträgheitsmoment die Trägheit gegenüber Änderungen der Winkelgeschwindigkeit.\TrueFalseAnswer{wahr}
\item Alle Gesetze der Drehbewegung sind formgleich zu entsprechenden Gesetzen der Translation.\TrueFalseAnswer{falsch}
\item Der Drehimpuls bezüglich des Kraftzentrums bleibt in einem Zentralkraftfeld erhalten.\TrueFalseAnswer{wahr}
\item Die kinetische Energie einer Rotationsbewegung ist proportional zur Winkelgeschwindigkeit $\omega$.\TrueFalseAnswer{falsch}
\item Zieht eine Eiskunstläuferin während einer Pirouette die Arme an sich heran, reduziert sich ihr (Massen-)Trägheitsmoment.\TrueFalseAnswer{wahr}
\item Zieht eine Eiskunstläuferin während einer Pirouette die Arme an, wird ihre Winkelgeschwindigkeit kleiner.\TrueFalseAnswer{falsch}
\item Zieht eine Eiskunstläuferin bei vernachlässigbarem äußerem Drehmoment die Arme an, verändert sich ihr Drehimpuls.\TrueFalseAnswer{falsch}
\item Zieht eine Eiskunstläuferin während einer Pirouette die Arme an, ändert sich ihre Rotationsenergie.\TrueFalseAnswer{wahr}
\item Präzession tritt auf, wenn ein Drehmoment $\vec{M}$ parallel zum Drehimpuls $\vec{L}$ einer Rotationsbewegung wirkt.\TrueFalseAnswer{falsch}
\end{enumerate}
\end{ExerciseBlock}
